\documentclass[11pt]{article}

\usepackage{graphicx}
\usepackage{amsmath}
\usepackage{booktabs}
\usepackage{hyperref}

\title{
From LeetCode to Prompt Engineering:
Investigating the Impact of LLM-Assisted Programming on Algorithmic Problem Solving
}

\author{
Chunru Qiu\\
Concordia University
}

\date{2026}

\begin{document}

\maketitle


\begin{abstract}

Large Language Models (LLMs) are transforming software engineering workflows.
Traditionally, computer science students develop algorithmic problem-solving
skills through extensive LeetCode practice. This research investigates whether
prompt engineering and AI-assisted programming can change the traditional
algorithm training paradigm.

\end{abstract}


\section{Introduction}

The development of Large Language Models has introduced new approaches
for software development. Developers can now use AI systems to generate,
optimize, and debug code.

Traditional software engineering preparation relies heavily on data structures,
algorithms, and problem-solving practice.

This paper investigates whether optimized prompt engineering can improve
algorithmic problem-solving performance and reduce reliance on traditional
LeetCode-style preparation.


\section{Research Question}

The main research question is:

\textbf{
Can optimized prompt engineering improve algorithmic problem-solving
performance compared with traditional LeetCode-style approaches?
}


\section{Methodology}

The proposed methodology compares three prompting strategies:

\begin{itemize}
    \item Baseline Prompt
    \item Optimized Algorithm Reasoning Prompt
    \item AI-Assisted Engineering Workflow
\end{itemize}


\section{Experimental Design}

The benchmark will contain LeetCode-style problems from different algorithm
categories.

Evaluation metrics include:

\begin{itemize}
    \item Correctness
    \item Time complexity
    \item Space complexity
    \item Runtime performance
\end{itemize}


\section{Expected Contribution}

This research aims to understand how AI-assisted programming may influence
future software engineering skills and education.


\section{Conclusion}

Future work will implement an automated benchmark system and evaluate
different prompt strategies using Large Language Models.

\end{document}
